\chapter{Implementierungsdetails}


\section{Web App}
\subsection{Die Zentrierung service-relevanter Einstellungen}
% extrahierung von theme colors etc.

\subsection{Zwischenspeicherstrategie im Service Worker}
Ein Service Worker ist ein in JavaScript geschriebener,
eventbasierter Proxy, welcher in einem separaten Thread im Browser
läuft und an eine spezifische Origin gekoppelt ist. Er hat
die Möglichkeit Anfragen zwischen der im Browser ausgeführten
Anwendung und dem Server abzufangen und zu verarbeiten.

Cache Only
Network Only
Cache First
Network First
Cache then Network
Stale while Revalidate
Generic Fallback

\begin{description}
    \item[Für den Zwischenspeicher des Browsers relevante Untergliederung]~\par
    \begin{itemize}
       \item Statische Dateien mit Hash
       \begin{itemize}
            \item main.dd5a1ad0.chunk.css
            \item main.46e36a4b.chunk.js
            \item 2.dc039c03.chunk.js
       \end{itemize}
       \item Statische Dateien ohne Hash
       \begin{itemize}
            \item index.html
            \item manifest.json
            \item favicon.ico
       \end{itemize}
       \item Antwort auf AJAX Anfragen in JSON
       \item Plugin Bundle Quellcode
    \end{itemize}
    \item[Zwischenspeicherstrategie]~\par
    \begin{enumerate}
       \item Statische Dateien mit durch Webpack generierten Hash im Dateinamen (, , )
       \begin{enumerate}
          \item erstens.
          \item zweitens.
       \end{enumerate}
       \item Statische Dateien ohne Hash im Dateinamen (index.html, manifest.json, favicon.ico)
    \end{enumerate}
 \end{description}

\begin{enumerate}
    \item Statische Dateien mit durch Webpack generierten Hash im Dateinamen
    \item static assets ohne hashes werden network first gecached
    \item AJAX calls werden network first, then indexedDB, then redirect to offline page
\end{enumerate}

\subsection{Seitengenerator}
% die verschiedenen Seiten sollten für das caching evtl. mit einem hash on build time versehen werden.


\subsection{Stetige Aktualisierung der PWA}
Probleme: cookies verändern sich, serviceworker verändert sich, neue Felder in der Datenbank
Frontend und Backend etc, also allen verschiedenen Services

\section{Benutzerverwaltung}
\subsection{Mailserver}
\subsection{JWT}
\subsection{Two-factor Auth}
\subsection{Two-Way Validation}
\subsection{Die Rechteverwaltung}
\subsection{Whitelist statt Blacklist}
\subsection{One verification route}
\subsection{Horizontale Skalierung ermöglichen}
% stateless backend, no sessions etc.


\section{Datenlieferung}

\section{Plugins}
\subsection{Code Splitting und Lazy Loading}
% laoding plugins on runtime
\subsection{D3.js und Chart.js}
\begin{listing}
    \label{lst:HelloJSX}
    \caption{Ein einfaches JSX Beispiel}
    \inputminted{jsx}{snippets/examples/Welcome.jsx}
\end{listing}

\begin{listing}
    \label{lst:Golang}
    \caption{Ein einfaches Golang Beispiel}
    \inputminted{go}{snippets/examples/hello.go}
\end{listing}

\newpage

\begin{figure}
    \label{figure:test}
    \includegraphics[scale=0.2]{img/HTW}
    \caption{Test}
\end{figure}

\begin{figure}
    \label{figure:beispiel}
    \includegraphics[scale=0.2]{img/HTW}
    \caption{Beispiel}
\end{figure}
